\section{Catálogo reproducible de gráficas}
\label{sec:catalogo-graficas}

Esta sección cubre todas las funciones públicas cuyo nombre comienza con
\codepath{plot_} o \codepath{render_} en
\codepath{hidden_attractors.plotting}: 33 funciones en la versión documentada.
Cada imagen fue producida llamando a la librería, no dibujada manualmente. Las
entradas tampoco son curvas pedagógicas ni matrices sintéticas construidas a
mano: proceden de cálculos numéricos reproducibles del sistema registrado
\codepath{chua-nonsmooth}. El caso
\codepath{chua_nonsmooth_real_system_catalog_20260729} es una reintegración
con la implementación canónica de la biblioteca.

El vector usado es
\[
(\alpha,\beta,\gamma,m_0,m_1,a_1,a_2,\rho)=
(8.4562,12.0732,0.0052,-0.1768,-1.1468,0.4,-1.5585,1).
\]
Según la salida, el contrato numérico corresponde a una trayectoria
\(q=1\) con \codepath{efork_q1}, una continuación entera en \(\lambda\), un
barrido de \(\beta\) con RK4, exponentes finito-temporales con QR--Benettin,
controles integrados desde vecindades de los equilibrios o una clasificación
de cuenca por tiempo finito. El manifiesto conserva para cada insumo el paso,
horizonte, condición inicial, malla, umbrales y semilla exactos.

Estos son resultados numéricos reales de un sistema matemático definido, no
mediciones experimentales ni una nueva campaña de validación. Ninguna imagen
certifica por sí sola caos, estabilidad asintótica, ocultedad, convergencia de
un integrador o desempeño científico.

El catálogo completo se regenera desde la raíz de \codepath{version_2} con:
\begin{verbatim}
python figure_scripts/generate_plot_catalog_examples.py
\end{verbatim}
El archivo
\codepath{docs/assets/generated_plot_catalog/catalog_results.json} conserva
la firma inspeccionada, el comando individual, la imagen representativa y el
estado de cada función, así como las 62 imágenes y sus hashes SHA-256. El
generador compara su lista con
\codepath{hidden_attractors.plotting.__all__} y falla si la API y el catálogo
dejan de coincidir.

\newcommand{\plotcatalogentry}[4]{%
  \bigskip
  \noindent\textbf{Función.}\ \codepath{#1}\par
  \noindent #2\par
  \noindent\textbf{Llamada directa.}\ \codepath{#3}\par
  \noindent\textbf{Comando del ejemplo.}\
  \texttt{python}\ \codepath{figure_scripts/generate_plot_catalog_examples.py}\
  \texttt{--only}\ \codepath{#1}\par
  \begin{center}
    \includegraphics[
      width=0.82\textwidth,
      height=0.48\textheight,
      keepaspectratio
    ]{#4}\\[-1mm]
    {\footnotesize Sistema \codepath{chua-nonsmooth} con el vector de
    parámetros declarado arriba. La entrada procede de la reintegración
    canónica y del contrato específico registrado en
    \codepath{catalog_results.json}. Es una salida numérica reproducible de la
    API, no una medición ni nueva evidencia de validación; por sí sola no
    demuestra caos, estabilidad, ocultedad o convergencia.}
  \end{center}
}

\newcommand{\plotcatalogextra}[3]{%
  \medskip
  \noindent\textbf{#2} (\codepath{#1}).\par
  \begin{center}
    \includegraphics[
      width=0.76\textwidth,
      height=0.39\textheight,
      keepaspectratio
    ]{#3}\\[-1mm]
    {\footnotesize Salida adicional calculada para
    \codepath{chua-nonsmooth}; parámetros, contrato numérico, procedencia y
    hash se registran en \codepath{catalog_results.json}. No constituye una
    medición experimental, nueva validación ni prueba autónoma de una
    propiedad dinámica.}
  \end{center}
}

\subsection{Trayectorias, series y espectros}

\plotcatalogentry
  {plot_phase_space}
  {Representa una trayectoria en dos o tres coordenadas, con color temporal
  opcional. La entrada usual tiene columnas \((t,x,y,z)\).}
  {plot_phase_space(trajectory, output_path, dims=("x","y","z"))}
  {assets/generated_plot_catalog/examples/08_plot_phase_space.png}

\plotcatalogentry
  {plot_phase_projections}
  {Construye las proyecciones \(x\)-\(y\), \(x\)-\(z\) y \(y\)-\(z\) de una
  misma trayectoria.}
  {plot_phase_projections(trajectory, output_path)}
  {assets/generated_plot_catalog/examples/07_plot_phase_projections.png}

\plotcatalogentry
  {plot_time_series}
  {Grafica las componentes seleccionadas contra el tiempo de muestreo.}
  {plot_time_series(trajectory, output_path, columns=("x","y","z"))}
  {assets/generated_plot_catalog/examples/10_plot_time_series.png}

\plotcatalogentry
  {plot_trajectory_overlay}
  {Superpone dos o más trayectorias etiquetadas para una comparación visual
  bajo la misma proyección.}
  {plot_trajectory_overlay(trajectories, labels, title="Comparacion", output_path=output_path)}
  {assets/generated_plot_catalog/examples/12_plot_trajectory_overlay.png}

\plotcatalogentry
  {plot_spectrum}
  {Representa un \codepath{SpectrumResult} calculado previamente, con
  frecuencia en las unidades declaradas.}
  {plot_spectrum(spectrum, output_path, x_units="Hz")}
  {assets/generated_plot_catalog/examples/09_plot_spectrum.png}

\plotcatalogentry
  {plot_trajectory_spectra}
  {Calcula y exporta el espectro de las componentes de una trayectoria. Acepta
  FFT y los alias \codepath{psd}, \codepath{welch} y
  \codepath{psd_welch}.}
  {plot_trajectory_spectra(trajectory, output_dir, method="fft", prefix="case")}
  {assets/generated_plot_catalog/examples/11_plot_trajectory_spectra.png}

\plotcatalogentry
  {plot_lyapunov_convergence}
  {Muestra la evolución finito-temporal de los exponentes contenidos en un
  \codepath{LyapunovResult}; no estima los exponentes por sí sola.}
  {plot_lyapunov_convergence(result, output_path)}
  {assets/generated_plot_catalog/examples/03_plot_lyapunov_convergence.png}

\plotcatalogentry
  {plot_attractor_trajectories}
  {Genera vistas de una trayectoria y de los equilibrios suministrados usando
  el contrato de configuración del workflow.}
  {plot_attractor_trajectories(trajectory, equilibria, config, output_dir)}
  {assets/generated_plot_catalog/examples/24_plot_attractor_trajectories.png}

\plotcatalogentry
  {plot_flexible_attractor_and_projections}
  {Exporta una vista espacial y proyecciones configurables para un caso
  identificado.}
  {plot_flexible_attractor_and_projections(trajectory, equilibria, config, output_dir, "case")}
  {assets/generated_plot_catalog/examples/25_plot_flexible_attractor_and_projections.png}

\plotcatalogentry
  {plot_timeseries_data}
  {Exporta las series temporales configuradas y el CSV correspondiente con
  columnas \(t,x,y,z\).}
  {plot_timeseries_data(trajectory, config, output_dir, "case")}
  {assets/generated_plot_catalog/examples/26_plot_timeseries_data.png}

\subsection{Bifurcación, cuencas y estabilidad}

\plotcatalogentry
  {plot_bifurcation_diagram}
  {Representa puntos ya postprocesados contra el parámetro; no continúa ramas
  ni identifica automáticamente bifurcaciones.}
  {plot_bifurcation_diagram(points, output_path, parameter_label="p", observable_label="x")}
  {assets/generated_plot_catalog/examples/06_plot_bifurcation_diagram.png}

\plotcatalogentry
  {plot_basin_slices}
  {Representa una o varias matrices de clases sobre planos declarados. Los
  colores reflejan las etiquetas recibidas, no una prueba global de cuenca.}
  {plot_basin_slices(slices, "system_id", output_dir)}
  {assets/generated_plot_catalog/examples/13_plot_basin_slices.png}

\plotcatalogentry
  {plot_basin_slice_file}
  {Exporta un único corte de cuenca a partir de dos ejes y una matriz de
  clases.}
  {plot_basin_slice_file("xy", u, v, classes, "E0", "system_id", output_dir)}
  {assets/generated_plot_catalog/examples/14_plot_basin_slice_file.png}

\plotcatalogentry
  {plot_matignon_equilibria}
  {Ubica autovalores de los Jacobianos de los equilibrios respecto de la
  frontera angular de Matignon para el orden \(q\).}
  {plot_matignon_equilibria(system, equilibria, 0.9, output_dir)}
  {assets/generated_plot_catalog/examples/15_plot_matignon_equilibria.png}

\subsection{Transferencia de Lur'e y función descriptiva}

\plotcatalogentry
  {plot_lure_nyquist_describing_function}
  {Compara el lugar de Nyquist de \(W_q\) con la curva
  \(-1/N(A)\) y marca una semilla armónica suministrada.}
  {plot_lure_nyquist_describing_function(system.lure, seed, output_path, q=1.0)}
  {assets/generated_plot_catalog/examples/04_plot_lure_nyquist_describing_function.png}

\plotcatalogentry
  {plot_lure_transfer_components}
  {Representa componentes reales, imaginarias y magnitudes de la transferencia
  de Lur'e alrededor de una semilla.}
  {plot_lure_transfer_components(system.lure, seed, output_path, q=1.0)}
  {assets/generated_plot_catalog/examples/05_plot_lure_transfer_components.png}

\plotcatalogentry
  {plot_nyquist_transfer}
  {Representa valores complejos de transferencia sobre una malla de
  frecuencias y candidatos \((A,\omega,k)\).}
  {plot_nyquist_transfer(omega, W, candidates, config, output_dir)}
  {assets/generated_plot_catalog/examples/16_plot_nyquist_transfer.png}

\plotcatalogentry
  {plot_describing_function}
  {Representa la función descriptiva del sistema y los candidatos entregados
  por el usuario.}
  {plot_describing_function(system, candidates, config, output_dir)}
  {assets/generated_plot_catalog/examples/17_plot_describing_function.png}

\plotcatalogentry
  {plot_harmonic_residual_map}
  {Construye el mapa del residuo
  \(\lvert1+N(A)W(\mathrm{i}\omega)\rvert\). Un mínimo visual no reemplaza la
  verificación numérica de cierre.}
  {plot_harmonic_residual_map(system, candidates, config, output_dir)}
  {assets/generated_plot_catalog/examples/18_plot_harmonic_residual_map.png}

El ejemplo real de estas gráficas declara
\codepath{transfer_convention="opposite_sign"} y
\codepath{harmonic_condition="1_plus_WN"}; usa entonces
\[
W_{\mathrm{code}}(s)=c^\mathsf{T}(P-sI)^{-1}b,\qquad
\left|1+N(A)W_{\mathrm{code}}(\mathrm{i}\omega)\right|,
\qquad W_{\mathrm{code}}(\mathrm{i}\omega_0)=-\frac{1}{k}.
\]
La pareja normalizada también aceptada por la API es
\codepath{standard + 1_minus_WN}. Como
\(W_{\mathrm{report}}=-W_{\mathrm{code}}\), en ella se representa
\(\lvert1-N(A)W_{\mathrm{report}}\rvert\) y el cierre queda en
\(W_{\mathrm{report}}=1/k\). Las funciones infieren la clave complementaria
si se proporciona solamente una y rechazan parejas de signo incoherentes,
salvo habilitación explícita.

\subsection{Continuación y controles de vecindad}

\plotcatalogentry
  {plot_integer_lure_continuation}
  {Representa los pasos de una continuación de Lur'e de orden entero ya
  calculada.}
  {plot_integer_lure_continuation(steps, output_path)}
  {assets/generated_plot_catalog/examples/02_plot_integer_lure_continuation.png}

\plotcatalogentry
  {plot_integer_hiddenness_controls}
  {Compara la trayectoria objetivo con controles enteros suministrados desde
  vecindades de equilibrios.}
  {plot_integer_hiddenness_controls(trajectory, probes, output_path)}
  {assets/generated_plot_catalog/examples/01_plot_integer_hiddenness_controls.png}

\plotcatalogentry
  {plot_continuation_eta}
  {Exporta norma final y amplitud pico a pico contra el parámetro de
  continuación \(\lambda\), además de comparaciones cuando hay varios pasos.
  El nombre público y los archivos con sufijo \codepath{_eta} se conservan
  únicamente por compatibilidad.}
  {plot_continuation_eta(cont_steps, config, output_dir)}
  {assets/generated_plot_catalog/examples/19_plot_continuation_eta.png}

\plotcatalogentry
  {plot_continuation_first_last_comparison}
  {Compara geométricamente las trayectorias del primer y último paso
  disponibles.}
  {plot_continuation_first_last_comparison(cont_steps, config, output_dir)}
  {assets/generated_plot_catalog/examples/20_plot_continuation_first_last_comparison.png}

\plotcatalogentry
  {plot_continuation_timeseries_comparison}
  {Compara las series temporales del primer y último paso de la continuación.}
  {plot_continuation_timeseries_comparison(cont_steps, config, output_dir)}
  {assets/generated_plot_catalog/examples/21_plot_continuation_timeseries_comparison.png}

\plotcatalogentry
  {plot_continuation_progression}
  {Muestra la progresión de estados de entrada, salida y trayectorias a lo
  largo de los pasos suministrados.}
  {plot_continuation_progression(cont_steps, config, output_dir)}
  {assets/generated_plot_catalog/examples/22_plot_continuation_progression.png}

\plotcatalogentry
  {plot_continuation_tracking}
  {Exporta por separado la norma rastreada y el estado discreto de cada paso.}
  {plot_continuation_tracking(cont_steps, config, output_dir)}
  {assets/generated_plot_catalog/examples/23_plot_continuation_tracking.png}

\plotcatalogentry
  {plot_neighborhood_control_spheres}
  {Representa la trayectoria, equilibrios y procedencia de sondas sobre
  esferas de control declaradas.}
  {plot_neighborhood_control_spheres(trajectory, probes, equilibria, config, output_dir)}
  {assets/generated_plot_catalog/examples/27_plot_neighborhood_control_spheres.png}

\plotcatalogentry
  {plot_sphere_test_results}
  {Resume los registros de una prueba de esfera para un equilibrio y radio
  específicos.}
  {plot_sphere_test_results("E0", equilibrium, radius, probe_runs, output_dir)}
  {assets/generated_plot_catalog/examples/28_plot_sphere_test_results.png}

\subsection{Renderizadores unificados}

\plotcatalogentry
  {render_attractor}
  {Renderiza una trayectoria y sus equilibrios en el almacén de figuras
  configurado, junto con metadatos.}
  {render_attractor(trajectory, equilibria, config, run_id="case")}
  {assets/generated_plot_catalog/examples/29_render_attractor.png}

\plotcatalogentry
  {render_basin}
  {Renderiza una matriz de cuenca y sus ejes en PNG, PDF y metadatos JSON.}
  {render_basin(grid_x, grid_y, basin_grid, config, run_id="case")}
  {assets/generated_plot_catalog/examples/30_render_basin.png}

\plotcatalogentry
  {render_nyquist}
  {Renderiza transferencia, función descriptiva y candidatos complejos con
  metadatos de procedencia.}
  {render_nyquist(freqs, W, N, candidates, config, run_id="case")}
  {assets/generated_plot_catalog/examples/31_render_nyquist.png}

\plotcatalogentry
  {render_matignon}
  {Renderiza autovalores y frontera de estabilidad para el orden declarado.}
  {render_matignon(eigenvalues, q, config, run_id="case")}
  {assets/generated_plot_catalog/examples/32_render_matignon.png}

\plotcatalogentry
  {render_all_plots}
  {Orquesta cualquier subconjunto compatible de atractor, cuenca, Nyquist y
  Matignon; no define un quinto diagnóstico matemático.}
  {render_all_plots(trajectory=trajectory, equilibria=equilibria, config=config, run_id="case")}
  {assets/generated_plot_catalog/examples/33_render_all_plots.png}

\subsection{Registro visual completo de funciones multisalida}

Las fichas anteriores conservan la correspondencia uno a uno entre las 33
funciones, sus llamadas directas y sus comandos individuales. Diez funciones
producen más de una imagen. El mismo comando
\codepath{python figure_scripts/generate_plot_catalog_examples.py --only FUNCTION_NAME}
indicado en cada ficha regenera su imagen principal y las salidas
adicionales siguientes. Con ellas, el inventario completo contiene 62 PNG:
33 representativos y 29 adicionales.

\plotcatalogextra
  {plot_trajectory_spectra}
  {Espectro de la componente 1}
  {assets/generated_plot_catalog/examples/11_plot_trajectory_spectra__02_catalog_fft_component_1.png}

\plotcatalogextra
  {plot_trajectory_spectra}
  {Espectro de la componente 2}
  {assets/generated_plot_catalog/examples/11_plot_trajectory_spectra__03_catalog_fft_component_2.png}

\plotcatalogextra
  {plot_nyquist_transfer}
  {Ampliación del entorno de cierre de Nyquist}
  {assets/generated_plot_catalog/examples/16_plot_nyquist_transfer__02_fig01b_nyquist_zoom_x.png}

\plotcatalogextra
  {plot_nyquist_transfer}
  {Partes real e imaginaria de la transferencia}
  {assets/generated_plot_catalog/examples/16_plot_nyquist_transfer__03_transfer_real_imag.png}

\plotcatalogextra
  {plot_continuation_eta}
  {Amplitud pico a pico contra \(\lambda\)}
  {assets/generated_plot_catalog/examples/19_plot_continuation_eta__02_continuation_amplitude_vs_eta.png}

\plotcatalogextra
  {plot_continuation_eta}
  {Comparación espacial del primer y último paso}
  {assets/generated_plot_catalog/examples/19_plot_continuation_eta__03_continuation_first_last_comparison.png}

\plotcatalogextra
  {plot_continuation_eta}
  {Proyecciones del primer y último paso}
  {assets/generated_plot_catalog/examples/19_plot_continuation_eta__04_continuation_first_last_projections.png}

\plotcatalogextra
  {plot_continuation_eta}
  {Comparación de series temporales}
  {assets/generated_plot_catalog/examples/19_plot_continuation_eta__05_continuation_timeseries_comparison_x.png}

\plotcatalogextra
  {plot_continuation_eta}
  {Progresión de la continuación}
  {assets/generated_plot_catalog/examples/19_plot_continuation_eta__06_continuation_progression.png}

\plotcatalogextra
  {plot_continuation_first_last_comparison}
  {Proyecciones del primer y último paso}
  {assets/generated_plot_catalog/examples/20_plot_continuation_first_last_comparison__02_continuation_first_last_projections.png}

\plotcatalogextra
  {plot_continuation_tracking}
  {Estado discreto por paso}
  {assets/generated_plot_catalog/examples/23_plot_continuation_tracking__02_continuation_tracking_status.png}

\plotcatalogextra
  {plot_attractor_trajectories}
  {Proyección \(xy\) de la trayectoria}
  {assets/generated_plot_catalog/examples/24_plot_attractor_trajectories__02_attractor_xy.png}

\plotcatalogextra
  {plot_attractor_trajectories}
  {Proyección \(xz\) de la trayectoria}
  {assets/generated_plot_catalog/examples/24_plot_attractor_trajectories__03_attractor_xz.png}

\plotcatalogextra
  {plot_attractor_trajectories}
  {Proyección \(yz\) de la trayectoria}
  {assets/generated_plot_catalog/examples/24_plot_attractor_trajectories__04_attractor_yz.png}

\plotcatalogextra
  {plot_flexible_attractor_and_projections}
  {Vista flexible en el plano \(xy\)}
  {assets/generated_plot_catalog/examples/25_plot_flexible_attractor_and_projections__02_catalog_flexible_xy.png}

\plotcatalogextra
  {plot_flexible_attractor_and_projections}
  {Vista flexible en el plano \(xz\)}
  {assets/generated_plot_catalog/examples/25_plot_flexible_attractor_and_projections__03_catalog_flexible_xz.png}

\plotcatalogextra
  {plot_flexible_attractor_and_projections}
  {Vista flexible en el plano \(yz\)}
  {assets/generated_plot_catalog/examples/25_plot_flexible_attractor_and_projections__04_catalog_flexible_yz.png}

\plotcatalogextra
  {plot_timeseries_data}
  {Serie temporal de la componente \(y\)}
  {assets/generated_plot_catalog/examples/26_plot_timeseries_data__02_catalog_series_timeseries_y.png}

\plotcatalogextra
  {plot_timeseries_data}
  {Serie temporal de la componente \(z\)}
  {assets/generated_plot_catalog/examples/26_plot_timeseries_data__03_catalog_series_timeseries_z.png}

\plotcatalogextra
  {plot_timeseries_data}
  {Series temporales conjuntas \(x,y,z\)}
  {assets/generated_plot_catalog/examples/26_plot_timeseries_data__04_catalog_series_timeseries_xyz.png}

\plotcatalogextra
  {render_attractor}
  {Renderizado del atractor en el plano \(xy\)}
  {assets/generated_plot_catalog/examples/29_render_attractor__02_chua-nonsmooth_attractor_xy.png}

\plotcatalogextra
  {render_attractor}
  {Renderizado del atractor en el plano \(xz\)}
  {assets/generated_plot_catalog/examples/29_render_attractor__03_chua-nonsmooth_attractor_xz.png}

\plotcatalogextra
  {render_attractor}
  {Renderizado del atractor en el plano \(yz\)}
  {assets/generated_plot_catalog/examples/29_render_attractor__04_chua-nonsmooth_attractor_yz.png}

\plotcatalogextra
  {render_all_plots}
  {Renderizado conjunto del atractor en \(xy\)}
  {assets/generated_plot_catalog/examples/33_render_all_plots__02_chua-nonsmooth_attractor_xy.png}

\plotcatalogextra
  {render_all_plots}
  {Renderizado conjunto del atractor en \(xz\)}
  {assets/generated_plot_catalog/examples/33_render_all_plots__03_chua-nonsmooth_attractor_xz.png}

\plotcatalogextra
  {render_all_plots}
  {Renderizado conjunto del atractor en \(yz\)}
  {assets/generated_plot_catalog/examples/33_render_all_plots__04_chua-nonsmooth_attractor_yz.png}

\plotcatalogextra
  {render_all_plots}
  {Renderizado conjunto de la cuenca finito-temporal}
  {assets/generated_plot_catalog/examples/33_render_all_plots__05_chua-nonsmooth_basin.png}

\plotcatalogextra
  {render_all_plots}
  {Renderizado conjunto de Nyquist}
  {assets/generated_plot_catalog/examples/33_render_all_plots__06_chua-nonsmooth_nyquist.png}

\plotcatalogextra
  {render_all_plots}
  {Renderizado conjunto de Matignon}
  {assets/generated_plot_catalog/examples/33_render_all_plots__07_chua-nonsmooth_matignon.png}

\noindent
El nombre \codepath{plot_continuation_eta} y los sufijos
\codepath{_eta} de algunas rutas se mantienen sólo por compatibilidad. Las
etiquetas visibles usan el parámetro efectivo \(\lambda\). El manifiesto
generado exige 33 resultados exitosos y enumera cada una de estas 62 imágenes
en \codepath{produced_outputs}.
